\chapter{氢原子}
\label{ch:hydrogen}

氢原子放电光谱由一系列离散谱线组成。
测量各谱线的波长后，Balmer型整数关系把看似无关的频率组织成不同系列。
同一系列共享一个较低能级，较高能级逐步接近电离阈值。
电子与原子核之间的 Coulomb 相互作用为何给出与 $1/n^2$ 成比例的束缚能量？具有不同角动量的态又为何可以同能？

观测到的光谱频率只给能量差。
取无限远分离的电子与原子核为内部能量零点后，
束缚态能量为负，谱线终点与电离阈值才共同固定能级的绝对标度。
此外原子核具有有限质量；同一套电子相互作用在不同同位素中
应产生可预测的小位移，这提供了两体模型的独立检验。

\section{两体运动的约化}
角动量代数给出了同一多重态的取向数。求谱线的波长还需要确定能级差。
对氢原子，束缚来自电子和原子核之间的静电吸引。分离质心动能后，
两者的相对运动决定原子的内部能级。

把原子核视为电荷$Ze$、质量$M$的点粒子，电子质量为$m_\electron$。
忽略自旋、相对论和辐射修正，记约化质量(reduced mass)为$\mu$、Coulomb耦合常数为$\kappa$：
\begin{equation}
    \mu = \frac{m_\electron M}{m_\electron + M},\qquad
    \kappa = \frac{Ze^2}{4\pi\varepsilon_0}.
\end{equation}

电子和核的位置算符与动量算符分别记为
$(\opbm{X}_\electron,\opbm{P}_\electron)$与$(\opbm{X}_\nucleus,\opbm{P}_\nucleus)$。
定义质心(center of mass)和相对变量
\begin{subequations}\label{eq:cm-relative}
    \begin{alignat}{2}
        \opbm{X}_{\CM} &= \frac{m_\electron\opbm{X}_\electron + M\opbm{X}_\nucleus}{m_\electron + M},
        &\quad \opbm{X} &= \opbm{X}_\electron - \opbm{X}_\nucleus,\\
        \opbm{P}_{\CM} &= \opbm{P}_\electron + \opbm{P}_\nucleus,
        &\quad \opbm{P} &= \frac{M\opbm{P}_\electron - m_\electron\opbm{P}_\nucleus}{m_\electron + M}.
    \end{alignat}
\end{subequations}
这些变量满足两组相互独立的正则对易关系。
动量的逆变换为
\begin{subequations}
    \begin{alignat}{1}
        \opbm{P}_\electron &= \frac{m_\electron}{m_\electron + M}\opbm{P}_{\CM} + \opbm{P},\\
        \opbm{P}_\nucleus &= \frac{M}{m_\electron + M}\opbm{P}_{\CM} - \opbm{P}.
    \end{alignat}
\end{subequations}
代入动能，两项交叉项分别为
$+\opbm{P}_{\CM} \cdot \opbm{P}/(m_\electron + M)$和它的负值，因而消去。

余下的平方项给出
\begin{subequations}
    \begin{align}
        \op{H}_\tot &= \frac{\opbm{P}_{\CM}^2}{2(m_\electron + M)} + \op{H},\\
        \op{H} &= \frac{\opbm{P}^2}{2\mu} - \kappa\op{R}^{-1},
        \quad \op{R} = \abs{\opbm{X}}.
    \end{align}
\end{subequations}
在位置表象中，距离算符 $\op R$ 使波函数乘以 $r = \abs{\bm r}$，Coulomb 势能项使波函数乘以 $-\kappa/r$。
自由质心的连续能谱与内部能级属于不同自由度。
后面$E$专指相对运动(relative motion)能量，零点取无穷远电离阈值(ionization threshold)。
点核模型只含 Coulomb 吸引势，径向解在原点取正则分支（\secref{sec:hydrogen-radial}）。

\section{氢样原子的基态波函数与特征半径}
\label{sec:hydrogen-1s}
相对运动的一个可观测量是电子到原子核的距离。若探测器只记录距离，球对称束缚态的读数将如何分布？
Coulomb 势具有球对称性，可以先求振幅只依赖距离的定态。
径向方程同时确定波函数的衰减长度和能量。

写这个球对称态为 $\psi(\bm r) = f(r)$。
位置表象中 $\opbm P^2 = -\hbar^2\laplacian$。对球对称的 $f$，
梯度沿径向，穿过半径 $r$ 球面的梯度通量是 $4\pi r^2 f'(r)$。
在薄壳的内、外表面取通量差，再除以壳体积 $4\pi r^2\d r$，
得 $\laplacian f = (r^2 f')'/r^2 = f'' + 2f'/r$。定态方程因此写作
\begin{equation}
    -\frac{\hbar^2}{2\mu}\biggp{f'' + \frac{2}{r}f'} - \frac{\kappa}{r}f = Ef.
\end{equation}
在远处势能趋于零，负能量解必须选随距离衰减的分支才能归一化。
试取处处有限、无节点的 $f = C\e^{-r/a}$，其中 $a > 0$ 待定。代回方程并除以 $f$，得
\[
    E = -\frac{\hbar^2}{2\mu a^2} +
        \frac{1}{r}\biggp{\frac{\hbar^2}{\mu a} - \kappa}.
\]
同一个定态在各处必须有相同的 $E$，所以 $1/r$ 的系数必须为零。
此条件使指数函数在所有 $r > 0$ 处都满足径向方程，同时确定了长度 $a_Z$ 和能量。
\begin{subequations}\label{eq:hydrogen-1s-scale}
    \begin{align}
        a_Z &= \frac{\hbar^2}{\mu\kappa},\\
        E_{1\mathrm{s}} &= -\frac{\hbar^2}{2\mu a_Z^2} = -\frac{\mu\kappa^2}{2\hbar^2}.
    \end{align}
\end{subequations}
其中$Z = 1, \mu = m_\electron$得到Bohr半径$a_0 = \hbar^2/(m_\electron e^2/4\pi\varepsilon_0)$。

取 $C$ 为正实数，球壳上的归一化条件给出
\[
    4\pi C^2\intRplus r^2\e^{-2r/a_Z}\d r = \pi a_Z^3 C^2 = 1 \implies C = \frac1{\sqrt{\pi a_Z^3}}.
\]
可得
\begin{equation}
    \psi_{1\mathrm{s}}(\bm r) = \frac{\e^{-r/a_Z}}{\sqrt{\pi a_Z^3}}.
\end{equation}
若探测器按半径等宽分箱，落入 $[r,r + \d r]$ 的概率为
$4\pi r^2\abs{\psi_{1\mathrm{s}}}^2\d r$。对 $r^2\e^{-2r/a_Z}$ 求导可得，计数最多的一箱在
$r = a_Z$ 附近，而放在等体积小格中的探测器在核附近计数最多。
这个态称为 $1\mathrm{s}$，其中正体字母 $\mathrm{s}$ 表示轨道角动量 $\ell = 0$。

\subsection{电子为什么不会落入原子核}
\label{ss:qm-atomic-stability}
若把电子当作绕固定核做圆周运动的经典点电荷，它有向心加速度，按经典电动力学就会辐射。
圆轨道的机械能为 $-\kappa/(2r)$，以辐射缓慢损失能量时，轨道半径会逐渐缩小。
在波动力学里，可用 $\psi_s(\bm r) = s^{3/2}\psi(s\bm r)$ 把一个波包的宽度压缩到原来的 $1/s$ 倍，其中 $s > 1$。
体积元缩小到 $1/s^3$ 倍，恰好抵消概率密度中的 $s^3$，总概率仍为一。
每求一次空间导数就多出一个因子 $s$，因此动能期望增至 $s^2$ 倍，Coulomb 吸引能的绝对值只增至 $s$ 倍。压缩得足够紧时，动能代价占主导。

允许波函数同时改变形状，还能找到比 $1\mathrm{s}$ 更低的能量吗？
记 $\bm n = \bm r/r$。展开下面的平方，交叉项含 $\bm n \cdot \grad\abs{\psi}^2$。
利用 $\div\bm n = 2/r$ 分部积分，并代入 $\hbar^2/(\mu a_Z) = \kappa$，对动能均值有限的归一化态得到
\begin{equation}
    \ave{\op H} - E_{1\mathrm{s}}
    = \frac{\hbar^2}{2\mu}\int_{\RR^3}
      \biggabs{\grad\psi + \frac{\bm n}{a_Z}\psi}^2\,\d^3 r \geq 0 .
    \label{eq:qm-hydrogen-energy-bound}
\end{equation}
$\psi_{1\mathrm{s}}$ 满足 $\grad\psi_{1\mathrm{s}} = -\bm n\psi_{1\mathrm{s}}/a_Z$，恰好达到下界。\footnote{分部积分可先取光滑且紧支撑的波函数。若从积分区挖去半径 $\epsilon$ 的小球，球面项随 $\epsilon^2$ 趋零。再用这类波函数逼近动能有限的态，即得同一下界。}
基态随时间只积累整体相位，其概率密度静止，径向概率流为零。耦合到辐射场后，激发态可跃迁到较低能级并发射光子。
孤立原子处于基态时，没有更低的电子能级可供这样的跃迁。自发发射的速率在\chapref{ch:transitions}由原子与场的联合态求出。

核本身能否束缚是另一个尺度上的问题。本章只把核的质量和电荷作为输入。
在由多个核子组成的核内，核子之间有短程核力，质子之间还有电斥力。氘核只有一个质子和一个中子，可以单独看出核力的束缚作用。
氘核的静能比自由质子和自由中子的静能之和低约 $\qty{2.22}{MeV}$，把它们拆开至少要补足这份结合能\cite{deuteron-binding}。
判断某个核是否会自发衰变，还要比较符合守恒律的各个末态的总能量。例如氚核束缚着一个质子和两个中子，仍可通过弱相互作用衰变为氦-3 核、电子和电子反中微子。
核的结合能和衰变能分别见《核辐射物理及探测学》的「原子核的结合能」和「$\beta$ 衰变」。

\section{Coulomb问题的守恒量}
球对称解给出了一个能量值，还没有说明不同角分布的激发态为什么可能同能。
普通中心势(central potential)使角动量守恒，旋转能联系同一 $\ell$ 中的不同投影态，但保持轨道角动量量子数 $\ell$。
要从对称性解释氢原子更高的简并度，需要寻找能够联系不同 $\ell$ 的守恒量。

这个守恒量可以从经典反平方力找出来。令 $\bm n = \bm r/r$，把速度分成径向和横向两部分，
$\bm v_\perp = \bm v - (\bm v \cdot \bm n)\bm n$，于是 $\dotbm{n} = \bm v_\perp/r$。
在力 $\dotbm{p} = -\kappa\bm n/r^2$ 下，$\bm L = \bm r \times \mu\bm v$ 守恒。
利用 $\bm n \times \bm L = -\mu r\bm v_\perp$ 可得
\[
    \dd{t}\frac{\bm p \times \bm L}{\mu}
       = -\frac{\kappa}{\mu r^2}\bm n \times \bm L
       = \kappa\dotbm{n}.
\]
两边作差就得到守恒的 $\bm p \times \bm L/\mu - \kappa\bm n$。
在位置表象中，径向单位矢量 $\bm n$ 对应的算符记作 $\opbm n$。其第 $i$ 个分量使波函数乘以径向与第 $i$ 根坐标轴夹角的余弦，
\[
    (\op n_i\psi)(\bm r) = \frac{r_i}{r}\,\psi(\bm r),\qquad r > 0.
\]
动量与角动量的乘积则要考虑 $\op{P}_j$ 与 $\op{L}_k$ 的作用次序，因此对经典乘积取对称排序。
\begin{definition}{Coulomb问题的Runge--Lenz向量}{qm-runge-lenz}
    在 Coulomb 势 $V(\op{R}) = -\kappa\op R^{-1}$ 中，相对运动的 Runge--Lenz 向量定义为
    \begin{equation}
        \opbm{A} = \frac{\opbm{P} \times \opbm{L} - \opbm{L} \times \opbm{P}}{2\mu}
        - \kappa\opbm n.
        \label{eq:RL}
    \end{equation}
    $\opbm{L} = \opbm{X} \times \opbm{P}$ 是相对轨道角动量。
    $\opbm{A}$ 的每个分量 $\op{A}_i$ 都具有能量乘长度的量纲，与 $\kappa$ 相同。
    其平方定义为 $\opbm{A}^2 = \opbm{A} \cdot \opbm{A}$。
\end{definition}
由$[\op{L}_j,\op{P}_k] = \i\hbar\epsilon_{jkl}\op{P}_l$，
\[
    \begin{alignedat}{1}
        (\opbm{L} \times \opbm{P})_i &= \epsilon_{ijk}\bigp{\op{P}_k\op{L}_j + [\op{L}_j,\op{P}_k]}\\
        &= -(\opbm{P} \times \opbm{L})_i + 2\i\hbar \op{P}_i.
    \end{alignedat}
\]
代回式\eqref{eq:RL}，就能读出对称排序产生的量子修正，
\begin{equation}
    \opbm A = \frac{\opbm P \times \opbm L - \i\hbar\opbm P}{\mu}
       - \kappa\opbm n.
\end{equation}
分子中正比于 $\hbar$ 的一项来自位置与动量的对易关系，各分量仍为对称算符。
先用只在某个球壳 $0 < r_1 < r < r_2$ 内非零的光滑波函数验算，由正则对易关系得到
\begin{subequations}
    \begin{align}
        [\op{H},\op{A}_i] &= 0,\\
        {}[\op{L}_i,\op{A}_j] &= \i\hbar\epsilon_{ijk}\op{A}_k,\\
        {}[\op{A}_i,\op{A}_j] &= -\frac{2\i\hbar \op{H}}{\mu}\epsilon_{ijk}\op{L}_k,\\
        \opbm{L} \cdot \opbm{A} &= \opbm{A} \cdot \opbm{L} = 0,\\
        \opbm{A}^2 &= \kappa^2\one + \frac{2\op{H}}{\mu}(\opbm{L}^2 + \hbar^2\one).
        \label{eq:RLsquare}
    \end{align}
\end{subequations}
为推导平方恒等式，令 $\opbm{Q} = \opbm{P} \times \opbm{L} - \i\hbar\opbm{P}$，
则 $\opbm{A} = \opbm{Q}/\mu - \kappa\opbm n$。其中 $\opbm n \cdot \opbm n = \one$，因为 $\textstyle \sum_i(r_i/r)^2 = 1$。另外两项满足
\begin{subequations}
    \begin{align}
        \opbm{Q}^2 &= \opbm{P}^2(\opbm{L}^2 + \hbar^2\one),\\
        \opbm{Q} \cdot \opbm n + \opbm n \cdot \opbm{Q}
        &= 2\op R^{-1}(\opbm{L}^2 + \hbar^2\one).
    \end{align}
\end{subequations}
将这两式代入$\opbm{A}^2$，逐步得到
\[
    \begin{alignedat}{1}
        \opbm{A}^2 &= \frac{\opbm{Q}^2}{\mu^2} + \kappa^2\one
            - \frac{\kappa}{\mu}
            \Bigp{\opbm{Q} \cdot \opbm n + \opbm n \cdot \opbm{Q}}\\
        &= \kappa^2\one +
            \Bigp{\frac{\opbm{P}^2}{\mu^2} - \frac{2\kappa}{\mu}\op R^{-1}}
            (\opbm{L}^2 + \hbar^2\one)
            = \kappa^2\one + \frac{2\op{H}}{\mu}(\opbm{L}^2 + \hbar^2\one).
    \end{alignedat}
\]
这就得到式\eqref{eq:RLsquare}，其中$\hbar^2$项来自量子排序。
这里$[\op{R},\opbm{L}^2] = [\opbm{P}^2,\opbm{L}^2] = 0$，因此可将$\opbm{L}^2 + \hbar^2\one$统一移到右侧。

其他对易式可逐分量用下面两式检验：
\begin{gather*}
    [\op{P}_i,f(\opbm{X})] = -\i\hbar(\p_i f)(\opbm{X}),\\
    \epsilon_{ijk}\epsilon_{imn} = \delta_{jm}\delta_{kn} - \delta_{jn}\delta_{km}.
\end{gather*}
点 Coulomb 势下的算符推导见文献\cite{runge}。

\subsectionstar{Runge--Lenz 矢量的守恒性}
\label{ss:rl-conservation}
守恒关系$[\op{H},\opbm{A}] = 0$可用两个相消的对易子直接验证。
对光滑的位置函数 $f$，由乘积求导得
\begin{equation}
    [\opbm{P}^2,f(\opbm{X})]
    = -2\i\hbar\sum_j(\p_j f)(\opbm X)\op{P}_j - \hbar^2(\laplacian f)(\opbm X).
\end{equation}
取普通函数 $n_i(\bm r) = r_i/r$，其导数为
$\p_j n_i = \delta_{ij}/r - r_i r_j/r^3$、$\laplacian n_i = -2r_i/r^3$（$r > 0$）。
分别代入位置算符 $\opbm X$，即得到上式所需的乘法算符。
把含动量的项合并为叉积，并用 $\op R^{-3}$ 表示乘以 $r^{-3}$ 的算符，得到
\begin{equation}
    \begin{alignedat}{1}
        \opbm{X} \times \opbm{L}
            &= \opbm{X}(\opbm{X} \cdot \opbm{P}) - \op{R}^2\opbm{P},\\
        [\opbm{P}^2,\opbm{n}]
            &= 2\op R^{-3}\bigp{\i\hbar\,\opbm{X} \times \opbm{L} + \hbar^2\opbm{X}}.
    \end{alignedat}
\end{equation}
另一方面$[\opbm{P}^2,\opbm{Q}] = 0$，而
\begin{equation}
    [\op{R}^{-1},\opbm{Q}]
    = -\op R^{-3}\bigp{\i\hbar\,\opbm{X} \times \opbm{L} + \hbar^2\opbm{X}}.
\end{equation}
因此
\begin{equation}
    [\op{H},\opbm{A}] = -\frac{\kappa}{\mu}[\op{R}^{-1},\opbm{Q}]
            - \frac{\kappa}{2\mu}[\opbm{P}^2,\opbm{n}] = 0.
\end{equation}
动能与Coulomb势的两组贡献逐项相消，说明$\opbm{A}$确实把同能态联系起来。

\section{束缚谱的代数推导}
$\opbm A$ 和 $\opbm L$ 的各分量均保持能量本征子空间。
在能量为 $E$ 的子空间内，$[\op{A}_i,\op{A}_j]$ 中 $\epsilon_{ijk}\op{L}_k$ 的系数为常数 $-2\i\hbar E/\mu$。
对于束缚态，这个性质使求能谱的问题可以归结为两个角动量的合成。

在固定负能量$E < 0$的子空间，定义
\begin{subequations}
    \begin{align}
        \opbm{D} &= \sqrt{-\frac\mu{2E}}\opbm{A},\\
        \opbm{M}_\pm &= \frac{\opbm{L} \pm \opbm{D}}{2}.
    \end{align}
\end{subequations}
此处$\opbm{D}$有角动量量纲。

由于$[\op{H},\op{A}_i] = 0$，$\op{A}_i$把固定$E$子空间映到自身，
而 $\op{H}$ 在这个子空间中作用为 $E\one$。
由缩放定义，三组对易关系变为
\begin{subequations}
    \begin{alignat}{1}
        [\op{L}_i,\op{L}_j] &= \i\hbar\epsilon_{ijk}\op{L}_k,\\
        [\op{L}_i,\op{D}_j] &= \i\hbar\epsilon_{ijk}\op{D}_k,\\
        [\op{D}_i,\op{D}_j] &= \i\hbar\epsilon_{ijk}\op{L}_k.
    \end{alignat}
\end{subequations}
例如$[\op{M}_{+i},\op{M}_{+j}]$展开后有四项，
\[
    \begin{alignedat}{1}
        \relax[\op{M}_{+i},\op{M}_{+j}]
        &= \frac14\bigp{[\op{L}_i,\op{L}_j] + [\op{L}_i,\op{D}_j] + [\op{D}_i,\op{L}_j] + [\op{D}_i,\op{D}_j]}\\
        &= \frac{\i\hbar}{4}\epsilon_{ijk}(2\op{L}_k + 2\op{D}_k)
        = \i\hbar\epsilon_{ijk}\op{M}_{+k}.
    \end{alignedat}
\]
$[\op{M}_{+i},\op{M}_{-j}]$中的$\opbm{L}$项和$\opbm{D}$项分别抵消。

因此得到两组彼此对易的角动量。
\begin{subequations}
    \begin{align}
        [\op{M}_{\pm i},\op{M}_{\pm j}] &= \i\hbar\epsilon_{ijk}\op{M}_{\pm k},\\
        {}[\op{M}_{+i},\op{M}_{-j}] &= 0.
    \end{align}
\end{subequations}
由 $\opbm{L} \cdot \opbm{D} = \opbm{D} \cdot \opbm{L} = 0$，有 $\opbm{M}_+^2 = \opbm{M}_-^2$。

两组角动量量子数因而相同，记为 $j$。在量子数为 $j$ 的子空间内，
\begin{equation}
    \opbm{M}_\pm^2 = \frac14(\opbm{L}^2 + \opbm{D}^2)
    = \hbar^2j(j + 1)\one.
\end{equation}
另一方面，将$\opbm{A}^2$恒等式乘$-\mu/(2E)$，并与上式联立，
\begin{equation}
    \opbm{L}^2 + \opbm{D}^2 = \biggp{-\frac{\mu\kappa^2}{2E} - \hbar^2}\one
    = 4\hbar^2j(j + 1)\one.
\end{equation}
把正整数 $n = 2j + 1$ 称为主量子数(principal quantum number)，便得到束缚谱。
\begin{theorem}{Coulomb束缚能谱}{qm-coulomb-spectrum}
    在上述非相对论点核 Coulomb 模型中，
    内部束缚态的能量为
    \begin{equation}
        E_n = -\frac{\mu\kappa^2}{2\hbar^2n^2}
            = -\frac{\mu c^2(Z\alpha)^2}{2n^2},
        \qquad n = 1,2,\ldots.
        \label{eq:hydrogen}
    \end{equation}
    $\mu$为电子与核的约化质量；未计入自旋时，
    固定$n$下的角动量为$\ell = 0,\ldots,n - 1$。
\end{theorem}

整数$n$来自恒等式
$4j(j + 1) + 1 = (2j + 1)^2$：
\[
    -\frac{\mu\kappa^2}{2E}
    = \hbar^2\bigb{4j(j + 1) + 1}
    = \hbar^2n^2.
\]
每一组角动量具有$2j + 1 = n$个投影态，
两个彼此对易的角动量因此构成$n \times n$维子空间。
\secref{sec:hydrogen-radial}将直接构造径向波函数，说明对每组允许的 $n,\ell$，径向解只差一个常数因子。

\begin{corollary}[氢原子束缚能级的简并度]
固定主量子数 $n$，合成 $\opbm{L} = \opbm{M}_+ + \opbm{M}_-$ 给出 $\ell = 0,1,\ldots,n - 1$，
每个 $\ell$ 有 $2\ell + 1$ 个磁量子数 $m$。因此空间态的总数是
\begin{equation}
    \sum_{\ell=0}^{n-1}(2\ell + 1) = n^2.
\end{equation}
计入两个电子自旋投影且忽略自旋耦合时，总简并度为 $2n^2$。
\end{corollary}
\begin{figure}[htbp]
    \centering\tikzfigure{hydrogen-levels}
    \caption{非相对论氢原子的前四层束缚态。纵向按 $-1/n^2$ 排列，横向按轨道角动量分列。红、青箭头各举一个 Balmer 通道：$3\mathrm{p} \to 2\mathrm{s}$ 对应 $\mathrm H_\upalpha$，$4\mathrm{d} \to 2\mathrm{p}$ 对应 $\mathrm H_\upbeta$。}\label{fig:hydrogen-levels}
\end{figure}
标准 Coulomb 问题的束缚态(bound state)记为 $\ket{n,\ell,m}$。三个标签依次给出同一个轨道态的能量、轨道角动量平方和 $z$ 投影的量子数。
位置与动量的矩阵元可结合这些量子数和算符恒等式进一步求取。
$E \geq 0$的连续谱也不适用上述两个紧致角动量代数。

\section{Rydberg 公式与同位素位移}
光谱仪测量波长。选取一对允许跃迁的能级，将能量差换算为光子波长，就可以把预测与观测逐线比较。
把氢换成氘时，核电荷相同，约化质量却变了；比较两种原子的同一条谱线，就能单独检验两体约化留下的质量因子。

记发射前、后原子的主量子数分别为 $n_\initial,n_\final$，其中 $n_\final < n_\initial$。忽略光子反冲和高阶修正时，发射线的波长满足
\begin{equation}
    \frac1\lambda = R_MZ^2\biggp{\frac1{n_\final^2} - \frac1{n_\initial^2}},
\end{equation}
其中$R_M$表示核质量为$M$时的Rydberg常数
\begin{equation}
    R_M = \frac{\mu e^4}{8\varepsilon_0^2h^3c}.
\end{equation}
而考虑无限重核极限$\mu = m_\electron$时，得到$R_\infty = \qty{10973731.568157}{m^{-1}}$。

固定 $n_\final = 2$，图\ref{fig:hydrogen-levels}中越来越密的上能级给出图\ref{fig:balmer-lines}的 Balmer 系列。令 $n_\initial \to \infty$ 得到短波极限。

\begin{figure}[htbp]
    \centering\tikzfigure{balmer-lines}
    \caption{氢的 Balmer 发射线（真空波长）。数字标出初态 $n_\initial$，红线 $\mathrm H_\upalpha$ 和青线 $\mathrm H_\upbeta$ 对应图\ref{fig:hydrogen-levels}中的同色箭头；左侧虚线标出约 $\qty{365}{nm}$ 的系限。}\label{fig:balmer-lines}
\end{figure}
将各实测波数除以括号中的整数因子后，
若模型适用，应得到同一个$R_MZ^2$，这比逐条拟合独立能量更强。

\begin{example}{同位素谱线的质量依赖}{qm-isotope}
    两个氢样离子的核电荷数均为 $Z$，核质量分别为 $M_1,M_2$。忽略光子反冲和精细结构，求同一电子跃迁的频率比。以无限重核时的频率 $\nu_\infty$ 为参照，求同位素频移对 $m_\electron/M_i$ 的一阶近似，并判断哪条谱线的波长较短。
    \par\smallskip\textbf{解}\quad
    固定初末态后，跃迁频率正比于电子与原子核的约化质量，故
    \begin{equation}
        \frac{\nu(M_2)}{\nu(M_1)}
        = \frac{\mu(M_2)}{\mu(M_1)},
        \qquad
        \mu(M) = \frac{m_\electron}{1 + m_\electron/M}.
    \end{equation}
    在$M \gg m_\electron$时展开分母，
    \begin{equation}
        \frac{\nu(M_2) - \nu(M_1)}{\nu_\infty}
        \simeq m_\electron\biggp{\frac1{M_1} - \frac1{M_2}}.
    \end{equation}
    因而较重同位素的频率更高、波长更短。
    这个质量依赖由质心分离产生，不需要改变电子电荷或Coulomb耦合。
\end{example}

有限核质量通过约化质量改变 $R_M$，产生同位素位移(isotope shift)。
精密谱还包括吸收发射时的整体反冲、精细结构、超精细结构、核大小和辐射修正。
谱线强弱则需要\chapref{ch:transitions}的跃迁矩阵元。

\subsectionstar{旧量子论为什么算对氢谱，又在哪里失效}
\label{ss:qm-old-quantum}
单凭氢的谱线位置，能判断电子在原子里怎样运动吗？在波动力学建立以前，Bohr 已用圆轨道模型得到了 Rydberg 公式\cite{bohrlecture}。
设电子与核相距 $r$，相对运动的速率为 $v$。Coulomb 吸引力提供向心加速度，故 $\mu v^2/r = \kappa/r^2$。
在经典力学中，圆轨道的角动量大小 $L = \mu vr$ 可以连续取值。Bohr 假定原子可以在一组特定轨道上运动而不辐射，
允许的圆轨道满足 $L = n\hbar$，$n = 1,2,\ldots$。消去 $v$ 后，半径与能量分别为
\begin{equation}
    r_n = \frac{L^2}{\mu\kappa} = n^2a_Z,\qquad
    E_n = \frac{\mu v^2}{2} - \frac{\kappa}{r_n}
       = -\frac{\kappa}{2r_n} = -\frac{\mu\kappa^2}{2\hbar^2 n^2}.
\end{equation}
模型还假定原子从高能轨道跃迁到低能轨道时，辐射频率满足 $h\nu = E_{n_\initial} - E_{n_\final}$。
代入上式就得到本节的谱线波长，连同 $R_M$ 对约化质量的依赖。

Sommerfeld 为什么引入椭圆？同一个 Coulomb 力已经允许椭圆形的经典束缚轨道。
设椭圆的半长轴为 $a$、离心率为 $\varepsilon$，Kepler 问题给出 $E = -\kappa/(2a)$ 和 $L^2 = \mu\kappa a(1 - \varepsilon^2)$。
固定 $a$ 而减小 $L$，轨道就变得更扁，能量保持不变。Sommerfeld 对每个独立周期施加了量子条件\cite{sommerfeld1916}。
在轨道平面取极坐标，沿运动方向量取方位角 $\varphi$，其共轭动量就是 $L$。径向动量 $p_r = \mu\dot r$ 在近核点与远核点处各变号一次。两个条件分别为
\begin{subequations}\label{eq:qm-bohr-sommerfeld-actions}
    \begin{align}
        \oint L\,\d\varphi &= 2\pi L = kh, & k &= 1,2,\ldots,\\
        \oint p_r\,\d r &= n_rh, & n_r &= 0,1,\ldots .
    \end{align}
\end{subequations}
两个积分相加就是沿椭圆一周的动量积分。设周期为 $\mathcal T$，动能为 $T$。
利用 $\d(\bm r \cdot \bm p)/\d t = 2T + V$ 在一周内的均值为零，可得 $\ave T = -E$。
代入 Kepler 周期 $\mathcal T = 2\pi\sqrt{\mu a^3/\kappa}$，得到
\[
    (n_r + k)h = \oint\bm p \cdot \d\bm r
        = 2\ave T\mathcal T = 2\pi\sqrt{\mu\kappa a}.
\]
将总作用量记为 $I = (n_r + k)\hbar$，即上式一周动量积分除以 $2\pi$。消去半长轴便得到
\begin{equation}\label{eq:qm-kepler-action-energy}
    E = -\frac{\mu\kappa^2}{2I^2}.
\end{equation}
这说明了旧量子论算对氢谱的原因。Kepler 轨道的能量只依赖总作用量，给定 $n = n_r + k$，所有形状都得到同一个 $E_n$。
再从波动方程来看，\secref{sec:hydrogen-radial}将说明，径向波函数既在原点正则又在无穷远衰减，要求
\begin{equation}\label{eq:qm-coulomb-radial-quantization}
    \frac{\kappa}{\hbar}\sqrt{\frac{\mu}{-2E}} = q + \ell + 1,\qquad q = 0,1,\ldots .
\end{equation}
这里 $q$ 是径向节点数，$\ell$ 是轨道角动量量子数。左端恰好就是上述 $I(E)/\hbar$，右端也取正整数。
两种方法因而选出完全相同的非相对论 Coulomb 能量，连 $n = 1$ 也相合。这个结论用到了 Coulomb 势的具体形式。

检验轨道形状，要看能量以外的读数。以 $n = 2$ 为例，旧理论允许 $k = 1$ 的椭圆和 $k = 2$ 的圆，
所以两种轨道的 $L^2/\hbar^2$ 分别为 $1,4$。波动力学允许 $\ell = 0,1$，相应的读数 $\ell(\ell + 1)$ 却是 $0,2$。
在基态，两种预言甚至可以用半径直接比较。Bohr 圆轨道给出确定的半径 $a_Z$ 和角动量 $\hbar$。
\secref{sec:hydrogen-1s}的 $1\mathrm{s}$ 态是球对称的，轨道角动量为零。对其径向概率密度积分，测得的核与电子距离 $R$ 小于 $a_Z$ 的概率为
\[
    \P(R < a_Z) = \int_0^{a_Z}\frac{4r^2}{a_Z^3}\e^{-2r/a_Z}\,\d r
             = 1 - 5\e^{-2} \simeq 0.32 .
\]
固定半径的轨道给这个事件的概率是零。波动力学中的 $a_Z$ 是径向概率密度的峰值位置。

高分辨率下，一条氢谱线可以分辨成相距很近的几条线。为解释这种精细结构，Sommerfeld 又采用了狭义相对论的动能。
取固定点核，动能按动量 $p$ 展开，对 $p^2/(2m_\electron)$ 的首项修正为 $-p^4/(8m_\electron^3c^2)$。
在同一能量的经典轨道上，$p^2$ 的周期平均值相同。圆轨道的 $p^2$ 恒定，椭圆轨道的 $p^2$ 随半径起伏，故后者的 $p^4$ 平均值更大。
由此可见相对论修正能够区分轨道形状。重新施加两个量子条件后，同一个 $n$ 的能量开始依赖 $k$。
旧模型所得的最低阶能移与\chapref{ch:relativity}的式\eqref{eq:fine}在取 $k = j + 1/2$ 时相合。旧模型的 $k$ 标记经典轨道的角动量，
现代的 $j$ 标记电子轨道角动量与自旋合成的总角动量。在 Dirac 理论的低速展开中，相对论动能、自旋轨道耦合和 Darwin 项共同给出这一能移。

谱线的强弱又提供了一个独立检验。忽略精细结构时，从 $2\mathrm{s}$ 和 $2\mathrm{p}$ 回到 $1\mathrm{s}$ 的能量差相同，
两个过程都满足同一个单光子能量条件。要预言哪条谱线有强度，现代理论用波函数求辐射的电偶极矩阵元。
$2\mathrm{s}$ 与 $1\mathrm{s}$ 都具有偶宇称，电偶极算符在空间反演下变号，因此两态间的矩阵元为零。
$2\mathrm{p}$ 具有奇宇称，到 $1\mathrm{s}$ 有非零矩阵元（\tabref{tab:hydrogen-dipole-z}）。
相应地，在隔离良好、无外场混合的氢原子中，$2\mathrm{p}$ 能迅速发射一个光子，而亚稳的 $2\mathrm{s}$ 主要通过较慢的双光子过程回到基态。
\ssecref{ss:lamb-experiment}的原子束实验正是用这两种态不同的存活率作读出\cite{lamb}。

作用量也能从波函数的相位求出。对在光滑势阱中往返的粒子，在两个转折点匹配波函数得到的量子化条件\eqref{eq:qm-wkb-bound}可写成
\[
    \oint p\,\d q \simeq \bigp{n + 1/2}h,
    \qquad n = 0,1,\ldots.
\]
多出的半整数记录了波在两端转向时积累的相位。当 $n$ 很大时，它只占总作用量的一小部分，旧量子条件由此有了半经典适用范围。
氢基态能量的精确相合则已由上面的 Coulomb 径向条件单独检验。

\section{径向波函数与节点}
\label{sec:hydrogen-radial}
光谱给出了能量。若在同一能级制备不同的轨道角动量，测到的电子与核之间的距离有何不同？
中心势允许把方向振幅与径向振幅分开，写成 $\psi = R(r)Y_{\ell m}(\theta,\phi)$。
球谐函数在球面上已归一化，因此球壳概率是 $r^2\abs{R(r)}^2\d r$。
球坐标中的动能分为径向项和角向项 $\opbm L^2/(2\mu r^2)$。
代入式\eqref{eq:qm-angular-laplacian}的本征值 $\hbar^2\ell(\ell + 1)$，径向方程为
\begin{equation}
    -\frac{\hbar^2}{2\mu}
       \biggp{R'' + \frac{2}{r}R' - \frac{\ell(\ell + 1)}{r^2}R}
       - \frac{\kappa}{r}R = E R.
    \label{eq:hydrogen-radial}
\end{equation}
令 $u = rR$，则归一化条件成为 $\textstyle \intRplus\abs u^2\d r = 1$，且 $R'' + 2R'/r = u''/r$。
对 $E < 0$ 定义衰减波数 $k = \sqrt{-2\mu E}/\hbar$、无量纲半径 $x = 2kr$ 和 $\eta = 1/(ka_Z)$，方程化为
\begin{equation}
    \dv[2]{u}{x} + \biggp{-\frac14 + \frac{\eta}{x} - \frac{\ell(\ell + 1)}{x^2}}u = 0 .
\end{equation}
通常的 Coulomb 边界条件在原点取正则支 $u \propto x^{\ell+1}$，在无穷远处取衰减支。
提取这两端的幂次和指数，令 $u = x^{\ell+1}\e^{-x/2}v(x)$。代回方程，
\begin{equation}
    xv'' + (2\ell + 2 - x)v' + (\eta - \ell - 1)v = 0,
\end{equation}
再展开
\begin{subequations}
    \begin{align}
        v &= \sum_{j=0}^\infty c_j x^j,\\
        c_{j+1} &= \frac{j + \ell + 1 - \eta}{(j + 1)(j + 2\ell + 2)}c_j.
    \end{align}
\end{subequations}
若级数不截断，$v(x)$ 在远处的增长会压过 $\e^{-x/2}$，径向波函数就不能归一化。\footnote{由递推式，高阶系数最终同号，且 $(j + 1)c_{j+1}/c_j$ 趋于 $1$。因此相邻系数的绝对值之比从某阶起大于 $3/[4(j + 1)]$。逐项与 $\e^{3x/4}$ 的级数比较，尾和的绝对值至少按 $\e^{3x/4}$ 增长。}
故递推的分子须在某个非负整数 $j = q$ 处为零。此时 $\eta = q + \ell + 1 = n$，于是 $k = 1/(na_Z)$。
由 $E = -\hbar^2 k^2/(2\mu)$ 再次得到式\eqref{eq:hydrogen}；每个允许的 $n,\ell$ 都有一个这样的多项式解。

记 $s = 2\ell + 1$。所得的 $q$ 次解正比于广义 Laguerre 多项式(generalized Laguerre polynomial)
\begin{equation}
    L_q^{(s)}(x) = \frac{x^{-s}\e^x}{q!}\dd[q]{x}\bigp{\e^{-x}x^{q+s}}.
\end{equation}
固定 $s$ 时，不同次数的多项式在权 $x^s\e^{-x}$ 下正交，平方范数为 $h_q = (q + s)!/q!$。
径向归一化的积分中多一个 $x$ 因子。三项递推中 $xL_q^{(s)}$ 沿 $L_q^{(s)}$ 的系数是 $2q + s + 1$，
其余项与 $L_q^{(s)}$ 正交，故径向归一化用到的积分为
\begin{equation}\label{eq:qm-laguerre-normalization-integral}
    \intRplus x^{s+1}\e^{-x}\bigb{L_q^{(s)}(x)}^2\d x
       = (2q + s + 1)h_q .
\end{equation}
写成 $x = 2r/(na_Z)$，并取原点附近径向振幅为正，归一化径向波函数为
\begin{equation}
    R_{n\ell}(r) = \biggp{\frac{2}{na_Z}}^{3/2}
       \sqrt{\frac{(n - \ell - 1)!}{2n(n + \ell)!}}\,
       x^\ell\e^{-x/2}L_{n-\ell-1}^{(2\ell+1)}(x).
    \label{eq:qm-hydrogen-radial-state}
\end{equation}
广义 Laguerre 多项式的 $q$ 个根彼此不同且都在正半轴上，故径向节点数为 $n - \ell - 1$。
例如 $n = 1$、$\ell = 0$ 时 $L_0^{(1)} = 1$，乘上 $Y_{00} = 1/\sqrt{4\pi}$，得到前面的 $\psi_{1\mathrm{s}} \propto \e^{-r/a_Z}$。

同能的 $2\mathrm{s}$ 和 $2\mathrm{p}$ 态在径向与角向都有差别。记 $\rho = r/a_Z$，在 $n = 2$ 时恰有 $x = \rho$。
对 $2\mathrm{s}$ 态，$L_1^{(1)}(\rho) = 2 - \rho$，电子的径向振幅在 $r = 2a_Z$ 过零；
对 $2\mathrm{p}$ 态，$L_0^{(3)} = 1$，径向振幅在正半轴上没有零点。
从角向看，$2\mathrm{s}$ 的方向分布均匀，而 $2\mathrm{p},m = 0$ 的 $Y_{10} = \sqrt{3/(4\pi)}\cos\theta$ 在赤道过零。
球壳概率密度 $P_{n\ell}(r) = r^2\abs{R_{n\ell}(r)}^2$ 给出图\ref{fig:hydrogen-radial}中的三条曲线：
\begin{subequations}
    \begin{align}
        a_Z P_{1\mathrm{s}} &= 4\rho^2\e^{-2\rho},\\
        a_Z P_{2\mathrm{s}} &= \frac18\rho^2(2 - \rho)^2\e^{-\rho},\\
        a_Z P_{2\mathrm{p}} &= \frac1{24}\rho^4\e^{-\rho}.
    \end{align}
\end{subequations}
$1\mathrm{s}$ 态的体概率密度从原点向外按指数衰减。乘上球壳的面积 $4\pi r^2$ 后，径向概率密度在 $r = a_Z$ 达到峰值（图\ref{fig:hydrogen-radial}）。
\begin{figure}[htbp]
    \centering\includegraphics[width=.62\linewidth]{figures/hydrogen_radial.pdf}
    \caption{氢样原子的 $1\mathrm{s}$、$2\mathrm{s}$、$2\mathrm{p}$ 径向分布。两个横坐标之间的曲线下面积给出电子落在相应球壳内的概率。$2\mathrm{s}$ 曲线在 $r/a_Z = 2$ 处触及横轴，标出两个峰之间的径向节点。}
    \label{fig:hydrogen-radial}
\end{figure}

\subsection{能量与速度的尺度}
利用Hellmann--Feynman关系（见\thmref{thm:qm-hellmann-feynman}及其证明）对$\kappa$求导，
\begin{subequations}
    \begin{align}
        -\ave{\op{R}^{-1}} &= \pv{E_n}{\kappa}
        = -\frac{\mu\kappa}{\hbar^2n^2},\\
        \ave{\op{R}^{-1}} &= \frac1{a_Zn^2}.
    \end{align}
\end{subequations}
要分别求动能 $\op T = \opbm P^2/(2\mu)$ 与势能 $\op V = -\kappa\op R^{-1}$ 的均值，引入缩放生成元
\begin{equation}
    \op G = \frac12(\opbm X \cdot \opbm P + \opbm P \cdot \opbm X).
\end{equation}
由正则对易关系，
\begin{subequations}
    \begin{alignat}{2}
        [\op{G},\op{X}_i] &= -\i\hbar \op{X}_i,
        &\quad [\op{G},\op{P}_i] &= \i\hbar \op{P}_i,\\
        \frac{\i}{\hbar}[\op{H},\op{G}] &= 2\op{T} - \opbm{X} \cdot (\grad V)(\opbm X).
    \end{alignat}
\end{subequations}
定态中 $\ave{\dot{\op G}} = 0$。对 Coulomb 势，$rV'(r) = -V(r)$，于是位力定理(virial theorem)给出 $2\ave{\op T} + \ave{\op V} = 0$。
与 $E_n = \ave{\op T} + \ave{\op V}$ 联立，就得到能量在动能和势能之间的分配，
\begin{equation}
    \ave{\op T} = -E_n,\qquad \ave{\op V} = 2E_n .
\end{equation}

这里的 Coulomb 束缚态在原点正则、在无穷远指数衰减，分部积分的边界项为零。有硬墙时，位力关系还包含壁面压力的贡献。
由此典型速度$\sqrt{\ave{\opbm{P}^2}}/\mu = Z\alpha c/n$，
非相对论近似的自然控制参数便是$Z\alpha/n$，对于低能级通常要求$Z\alpha \ll 1$。

\begin{example}{中心势的能级简并}{qm-central-potential-degeneracy}
    忽略自旋，设粒子在中心势中运动。旋转对称性对给定 $\ell$ 的束缚能级有什么要求？氢原子同一主量子数 $n$ 的空间态为何可以具有更高的简并度？
    \par\smallskip\textbf{解}\quad
    旋转只联系同一 $\ell$ 下的 $2\ell + 1$ 个 $m$ 态。Coulomb 问题的 Runge--Lenz 矢量分量 $\op{A}_i$ 还联系不同 $\ell$，两个角动量代数使 $\ell = 0,\ldots,n - 1$ 的各多重态具有同一能量。将它们的维数相加，得到 $n^2$ 个空间态。
\end{example}
\reading{Coulomb额外对称性参照Runge--Lenz代数资料\cite{runge}；
    径向波函数的一般式见 \href{https://dlmf.nist.gov/18.39.E37}{NIST 数学函数手册式18.39.37}。
    空间波函数、非相对论能谱及其精细结构修正可对照MIT 8.06讲义2第2.1至2.4节\cite{mitfine}。下一章研究偏离精确可解模型时怎样给出受控预测。}
